\chapter{Suites numériques}
\textit{Une suite, c'est une liste \emph{ordonnée} de nombres : le premier terme,
le deuxième, le troisième\dots\ Beaucoup de situations concrètes se modélisent ainsi
— l'argent placé à la banque qui grossit chaque année, une population qui évolue, un
remboursement échelonné. Ce chapitre installe le vocabulaire des suites, deux grandes
familles (arithmétiques et géométriques) que tu rencontreras partout, et les formules
qui permettent de calculer un terme lointain ou une somme sans tout additionner.}

\section{Notion de suite numérique}

\begin{definition}[Suite numérique]
Une \textbf{suite numérique} $u$ est une fonction définie sur $\N$ (ou sur une partie
$\{n\in\N \mid n\ge n_0\}$) et à valeurs dans $\R$. À l'entier $n$ elle associe un réel
noté $u_n$ (et non $u(n)$), appelé \textbf{terme de rang $n$}. La suite tout entière se
note $(u_n)$ ou $(u_n)_{n\ge n_0}$.
\end{definition}

\begin{remarque}
$u_n$ est le terme \emph{de rang} $n$ ; $u_0$ est le \textbf{premier terme} (parfois
$u_1$ selon l'énoncé). Il ne faut pas confondre le \emph{rang} $n$ (la position, un
entier) et le \emph{terme} $u_n$ (la valeur, un réel).
\end{remarque}

\section{Modes de génération d'une suite}

Pour décrire une suite, il faut une règle qui produit ses termes. On rencontre deux
grands procédés.

\begin{definition}[Suite définie par une formule explicite]
La suite $(u_n)$ est donnée par une \textbf{formule explicite} lorsqu'il existe une
fonction $f$ telle que, pour tout rang $n$,
\[ u_n = f(n). \]
On calcule alors n'importe quel terme \emph{directement}, sans connaître les précédents.
\end{definition}

\begin{exemple}
Soit $u_n = 2n+1$ pour $n\in\N$. Alors $u_0=1$, $u_1=3$, $u_2=5$ et, sans détour,
$u_{50}=2\times 50+1=101$.
\end{exemple}

\begin{definition}[Suite définie par récurrence]
La suite $(u_n)$ est définie par \textbf{récurrence} lorsqu'on se donne son premier terme
et une relation qui exprime chaque terme à partir du précédent :
\[ u_{0} \text{ connu},\qquad u_{n+1}=g(u_n)\ \text{ pour tout } n. \]
Pour obtenir $u_n$, on est obligé de calculer \emph{de proche en proche} tous les termes
qui le précèdent.
\end{definition}

\begin{exemple}
Soit $(u_n)$ définie par $u_0=3$ et $u_{n+1}=2u_n-1$. On obtient
$u_1=2\times 3-1=5$, puis $u_2=2\times 5-1=9$, puis $u_3=2\times 9-1=17$, etc.
\end{exemple}

\begin{remarque}
Une suite peut admettre les deux modes à la fois. Tout l'intérêt des suites
arithmétiques et géométriques (sections suivantes) est précisément de transformer une
récurrence en formule explicite, donc de calculer un terme lointain sans itérer.
\end{remarque}

\begin{illus}
\begin{tikzpicture}
\begin{axis}[
  width=11cm,height=6cm,
  axis lines=middle,
  xlabel={$n$},ylabel={$u_n$},
  xlabel style={right},ylabel style={above},
  xmin=-0.4,xmax=7.4,ymin=-0.5,ymax=8.5,
  xtick={1,2,3,4,5,6,7},ytick={2,4,6,8},
  tick style={onbmuted},
  every axis x label/.style={at={(ticklabel* cs:1)},anchor=west},
  every axis y label/.style={at={(ticklabel* cs:1)},anchor=south},
  clip=true]
  \addplot[only marks,onbo,mark=*,mark size=2pt] coordinates {
    (0,1)(1,3)(2,5)(3,7)(4,9)};
  \node[onbo] at (axis cs:4.7,8.0){$u_n=2n+1$};
\end{axis}
\end{tikzpicture}
\fig{Figure — nuage de points d'une suite : un point $(n\,;\,u_n)$ par rang. Les points
ne sont jamais reliés (le rang $n$ est un entier).}
\end{illus}

\section{Sens de variation d'une suite}

\begin{definition}[Suite croissante, décroissante]
Soit $(u_n)$ une suite. On dit qu'elle est :
\begin{itemize}
  \item \textbf{croissante} si $u_{n+1}\ge u_n$ pour tout $n$ (strictement croissante
        si $u_{n+1}>u_n$) ;
  \item \textbf{décroissante} si $u_{n+1}\le u_n$ pour tout $n$ (strictement
        décroissante si $u_{n+1}<u_n$) ;
  \item \textbf{monotone} si elle est croissante ou décroissante ;
  \item \textbf{constante} si $u_{n+1}=u_n$ pour tout $n$.
\end{itemize}
\end{definition}

\begin{aretenir}[Comment étudier le sens de variation]
On dispose de deux méthodes selon la forme de la suite.
\begin{itemize}
  \item \textbf{Signe de la différence} $u_{n+1}-u_n$ (méthode universelle) :
        si elle est $\ge 0$ pour tout $n$, la suite croît ; si elle est $\le 0$, elle
        décroît.
  \item \textbf{Comparaison du quotient à $1$}, valable \emph{uniquement si tous les
        termes sont strictement positifs} : on compare $\dfrac{u_{n+1}}{u_n}$ à $1$
        ($>1\Rightarrow$ croissante, $<1\Rightarrow$ décroissante).
\end{itemize}
\end{aretenir}

\begin{exemple}
Étudions $u_n=n^2-3n$. Pour tout $n$,
\[ u_{n+1}-u_n=\big((n+1)^2-3(n+1)\big)-\big(n^2-3n\big)=2n-2. \]
Ce nombre est $<0$ pour $n=0$ et $\ge 0$ dès que $n\ge 1$. La suite \emph{décroît}
de $u_0$ à $u_1$, puis \emph{croît} ensuite : elle n'est pas monotone sur $\N$ tout entier.
\end{exemple}

\section{Suites majorées, minorées, bornées}

\begin{definition}[Majorant, minorant, suite bornée]
Soit $(u_n)$ une suite et $M,m\in\R$.
\begin{itemize}
  \item $(u_n)$ est \textbf{majorée} par $M$ si $u_n\le M$ pour tout $n$ ;
  \item $(u_n)$ est \textbf{minorée} par $m$ si $u_n\ge m$ pour tout $n$ ;
  \item $(u_n)$ est \textbf{bornée} si elle est à la fois majorée et minorée,
        c'est-à-dire s'il existe $M\ge 0$ tel que $|u_n|\le M$ pour tout $n$.
\end{itemize}
\end{definition}

\begin{exemple}
La suite $u_n=\dfrac{1}{n+1}$ vérifie $0<u_n\le 1$ pour tout $n\in\N$ : elle est minorée
par $0$ et majorée par $1$, donc \textbf{bornée}.
\end{exemple}

\begin{remarque}
Un majorant (ou un minorant) n'est jamais unique : si $M$ majore la suite, tout réel
plus grand la majore aussi. On cherche en général le « meilleur » majorant, mais
n'importe lequel suffit pour prouver qu'une suite est majorée.
\end{remarque}

\section{Suites arithmétiques}

\begin{definition}[Suite arithmétique]
Une suite $(u_n)$ est \textbf{arithmétique} si l'on passe d'un terme au suivant en
ajoutant toujours le \emph{même} réel $r$, appelé \textbf{raison} :
\[ u_{n+1}=u_n+r \qquad \text{pour tout } n. \]
\end{definition}

\begin{aretenir}[Reconnaître une suite arithmétique]
$(u_n)$ est arithmétique $\iff$ la différence $u_{n+1}-u_n$ est \textbf{constante} ;
cette constante est alors la raison $r$. C'est le test à appliquer en premier.
\end{aretenir}

\begin{theoreme}[Terme général d'une suite arithmétique]
Si $(u_n)$ est arithmétique de raison $r$, alors pour tous rangs $n$ et $p$ :
\[ \boxed{\,u_n = u_p + (n-p)\,r\,}\qquad\text{en particulier}\qquad u_n=u_0+n\,r. \]
\end{theoreme}

\begin{exemple}[Démonstration]
En ajoutant $r$ à chaque étape de $u_p$ jusqu'à $u_n$, on franchit $n-p$ pas :
\[ u_n=\underbrace{u_p+\overbrace{r+r+\dots+r}^{n-p \text{ fois}}}=u_p+(n-p)\,r. \]
\end{exemple}

\begin{exemple}
$(u_n)$ est arithmétique, $u_3=7$ et de raison $r=2$. Alors
$u_{20}=u_3+(20-3)\times 2=7+34=41$, sans calculer les termes intermédiaires.
\end{exemple}

\begin{theoreme}[Somme de termes consécutifs]
Pour une suite arithmétique, la somme de termes consécutifs vaut le nombre de termes
multiplié par la \emph{moyenne} des termes extrêmes :
\[ \boxed{\,u_0+u_1+\dots+u_n=(n+1)\,\dfrac{u_0+u_n}{2}\,}. \]
En particulier, pour $r=1$ et $u_0=0$ : $\,0+1+2+\dots+n=\dfrac{n(n+1)}{2}$.
\end{theoreme}

\begin{exemple}[Démonstration]
Notons $S=u_0+u_1+\dots+u_n$. Écrivons cette somme dans les deux sens et additionnons
terme à terme :
\[
\begin{array}{rcccccccl}
S &=& u_0 &+& u_1 &+&\dots&+& u_n\\
S &=& u_n &+& u_{n-1} &+&\dots&+& u_0
\end{array}
\]
Chaque colonne $u_k+u_{n-k}$ vaut $u_0+u_n$ (car $u_k+u_{n-k}=2u_0+(k+n-k)r=u_0+u_n$).
Comme il y a $n+1$ colonnes, $2S=(n+1)(u_0+u_n)$, d'où le résultat.
\end{exemple}

\begin{illus}
\begin{tikzpicture}
\begin{axis}[
  width=11cm,height=6cm,
  axis lines=middle,
  xlabel={$n$},ylabel={$u_n$},
  xlabel style={right},ylabel style={above},
  xmin=-0.4,xmax=7.4,ymin=-0.5,ymax=8.5,
  xtick={1,2,3,4,5,6,7},ytick={2,4,6,8},
  tick style={onbmuted},
  every axis x label/.style={at={(ticklabel* cs:1)},anchor=west},
  every axis y label/.style={at={(ticklabel* cs:1)},anchor=south},
  clip=true]
  \addplot[onbmuted,dashed,thick,domain=0:7]{1+x};
  \addplot[only marks,onbo,mark=*,mark size=2pt] coordinates {
    (0,1)(1,2)(2,3)(3,4)(4,5)(5,6)(6,7)};
  \node[onbo] at (axis cs:5.4,4.6){$u_n=u_0+nr$};
\end{axis}
\end{tikzpicture}
\fig{Figure — les points d'une suite arithmétique ($u_0=1$, $r=1$) sont \emph{alignés} :
ils se trouvent sur une droite de pente $r$.}
\end{illus}

\section{Suites géométriques}

\begin{definition}[Suite géométrique]
Une suite $(u_n)$ est \textbf{géométrique} si l'on passe d'un terme au suivant en
multipliant toujours par le \emph{même} réel $q\neq 0$, appelé \textbf{raison} :
\[ u_{n+1}=q\,u_n \qquad \text{pour tout } n. \]
\end{definition}

\begin{aretenir}[Reconnaître une suite géométrique]
Lorsque aucun terme n'est nul, $(u_n)$ est géométrique $\iff$ le quotient
$\dfrac{u_{n+1}}{u_n}$ est \textbf{constant} ; cette constante est la raison $q$.
\end{aretenir}

\begin{theoreme}[Terme général d'une suite géométrique]
Si $(u_n)$ est géométrique de raison $q$, alors pour tous rangs $n$ et $p$ :
\[ \boxed{\,u_n = u_p\,q^{\,n-p}\,}\qquad\text{en particulier}\qquad u_n=u_0\,q^{\,n}. \]
\end{theoreme}

\begin{exemple}[Démonstration]
En multipliant par $q$ à chaque étape de $u_p$ jusqu'à $u_n$, on multiplie $n-p$ fois :
\[ u_n=u_p\times\underbrace{q\times q\times\dots\times q}_{n-p \text{ fois}}=u_p\,q^{\,n-p}. \]
\end{exemple}

\begin{exemple}
$(u_n)$ est géométrique, $u_0=5$ et $q=2$. Alors $u_6=5\times 2^6=5\times 64=320$.
\end{exemple}

\begin{theoreme}[Somme de termes consécutifs]
Si $(u_n)$ est géométrique de raison $q\neq 1$, alors
\[ \boxed{\,u_0+u_1+\dots+u_n=u_0\,\dfrac{1-q^{\,n+1}}{1-q}\,}. \]
(Si $q=1$, la suite est constante et la somme vaut simplement $(n+1)\,u_0$.)
\end{theoreme}

\begin{exemple}[Démonstration]
Posons $S=u_0+u_0q+u_0q^2+\dots+u_0q^{n}$. Multiplions par $q$ :
\[ qS=u_0q+u_0q^2+\dots+u_0q^{n}+u_0q^{n+1}. \]
En soustrayant, presque tout se télescope :
\[ S-qS=u_0-u_0q^{n+1}=u_0\,(1-q^{n+1}), \]
donc $S(1-q)=u_0(1-q^{n+1})$, d'où $S=u_0\,\dfrac{1-q^{n+1}}{1-q}$ car $q\neq 1$.
\end{exemple}

\begin{illus}
\begin{tikzpicture}
\begin{axis}[
  width=11cm,height=6cm,
  axis lines=middle,
  xlabel={$n$},ylabel={$u_n$},
  xlabel style={right},ylabel style={above},
  xmin=-0.4,xmax=6.4,ymin=-1,ymax=18,
  xtick={1,2,3,4,5,6},ytick={4,8,12,16},
  tick style={onbmuted},
  every axis x label/.style={at={(ticklabel* cs:1)},anchor=west},
  every axis y label/.style={at={(ticklabel* cs:1)},anchor=south},
  clip=true]
  \addplot[only marks,onbgreen,mark=*,mark size=2pt] coordinates {
    (0,1)(1,1.5)(2,2.25)(3,3.375)(4,5.0625)(5,7.59)(6,11.39)};
  \node[onbgreen] at (axis cs:4.1,13.5){$u_n=u_0\,q^{\,n}$};
\end{axis}
\end{tikzpicture}
\fig{Figure — pour $q>1$ ($u_0=1$, $q=1{,}5$), les points d'une suite géométrique
montent de plus en plus vite : la croissance est exponentielle.}
\end{illus}

\section{Applications : intérêts simples et composés}

Les suites arithmétiques et géométriques modélisent deux façons de faire fructifier un
capital, qu'on rencontre dans les banques et tontines au Cameroun.

\begin{propriete}[Intérêts simples — modèle arithmétique]
Un capital $C_0$ placé à \textbf{intérêts simples} au taux annuel $t$ rapporte chaque
année le \emph{même} intérêt $C_0\,t$. Le capital après $n$ années,
\[ C_n=C_0+n\,C_0\,t=C_0\,(1+n\,t), \]
forme une suite \textbf{arithmétique} de premier terme $C_0$ et de raison $r=C_0\,t$.
\end{propriete}

\begin{exemple}
On place $C_0=200\,000$ FCFA à intérêts simples au taux $t=5\%=0{,}05$. La raison vaut
$r=200\,000\times 0{,}05=10\,000$ FCFA par an. Après $4$ ans :
\[ C_4=200\,000\,(1+4\times 0{,}05)=200\,000\times 1{,}2=240\,000 \text{ FCFA}. \]
\end{exemple}

\begin{propriete}[Intérêts composés — modèle géométrique]
Un capital $C_0$ placé à \textbf{intérêts composés} au taux annuel $t$ voit, chaque
année, ses intérêts s'ajouter au capital, qui produit à son tour des intérêts. Le
capital après $n$ années,
\[ C_n=C_0\,(1+t)^{\,n}, \]
forme une suite \textbf{géométrique} de premier terme $C_0$ et de raison $q=1+t$.
\end{propriete}

\begin{exemple}
Mêmes données ($C_0=200\,000$ FCFA, $t=0{,}05$) mais à intérêts composés : $q=1{,}05$.
Après $4$ ans :
\[ C_4=200\,000\times 1{,}05^{\,4}\approx 200\,000\times 1{,}2155\approx 243\,101 \text{ FCFA}. \]
On dépasse les $240\,000$ FCFA des intérêts simples : les intérêts qui produisent à leur
tour des intérêts font la différence, d'autant plus nette que la durée est longue.
\end{exemple}

\begin{remarque}
Une \textbf{tontine} où chaque membre verse une part fixe et où le total grossit
régulièrement relève plutôt du modèle arithmétique ; un placement bancaire qui
\emph{capitalise} les intérêts relève du modèle géométrique. Savoir distinguer les deux,
c'est savoir choisir la bonne formule.
\end{remarque}

% ════════════════════════════════════════════════════════════
\section{L'essentiel à retenir}

\begin{synthese}
\textbf{Vocabulaire.} Une suite $(u_n)$ associe à chaque rang $n\in\N$ un réel $u_n$.
Génération : \emph{explicite} $u_n=f(n)$ (calcul direct) ou \emph{récurrente}
$u_{n+1}=g(u_n)$ avec $u_0$ donné (de proche en proche).

\smallskip
\textbf{Sens de variation.} Signe de $u_{n+1}-u_n$ (universel) ; ou, si $u_n>0$,
position de $\dfrac{u_{n+1}}{u_n}$ par rapport à $1$.

\smallskip
\textbf{Bornes.} Majorée : $u_n\le M$ ; minorée : $u_n\ge m$ ; bornée : les deux.

\smallskip
\textbf{Arithmétique} (raison $r$) : $u_{n+1}=u_n+r$, $\ u_n=u_p+(n-p)r$, et
\[ u_0+\dots+u_n=(n+1)\,\dfrac{u_0+u_n}{2}. \]
\textbf{Géométrique} (raison $q\neq 0$) : $u_{n+1}=q\,u_n$, $\ u_n=u_p\,q^{\,n-p}$, et
pour $q\neq 1$ :
\[ u_0+\dots+u_n=u_0\,\dfrac{1-q^{\,n+1}}{1-q}. \]

\smallskip
\textbf{Finance.} Intérêts simples $\to$ arithmétique : $C_n=C_0(1+nt)$.
Intérêts composés $\to$ géométrique : $C_n=C_0(1+t)^n$.
\end{synthese}

% ════════════════════════════════════════════════════════════
\section{Méthodes \& savoir-faire}

\methode{Calculer les premiers termes d'une suite récurrente}
Partir du premier terme donné, puis appliquer la relation $u_{n+1}=g(u_n)$ en cascade,
en réutilisant chaque résultat pour calculer le suivant.
\begin{exemple}
$u_0=1$ et $u_{n+1}=\dfrac{u_n+3}{2}$. On a $u_1=\dfrac{1+3}{2}=2$, $u_2=\dfrac{2+3}{2}=2{,}5$,
$u_3=\dfrac{2{,}5+3}{2}=2{,}75$.
\end{exemple}

\methode{Montrer qu'une suite est arithmétique (ou géométrique)}
Calculer $u_{n+1}-u_n$ : si c'est une \emph{constante}, la suite est arithmétique
(de raison cette constante). Sinon, si tous les termes sont non nuls, calculer
$\dfrac{u_{n+1}}{u_n}$ : si c'est constant, la suite est géométrique.
\begin{exemple}
$u_n=3\times 2^{\,n}$. Alors $\dfrac{u_{n+1}}{u_n}=\dfrac{3\times 2^{\,n+1}}{3\times 2^{\,n}}=2$,
constant : $(u_n)$ est géométrique de raison $q=2$ et de premier terme $u_0=3$.
\end{exemple}

\methode{Calculer un terme de rang élevé}
Utiliser la formule du terme général plutôt que d'itérer : $u_n=u_p+(n-p)r$ pour une
arithmétique, $u_n=u_p\,q^{\,n-p}$ pour une géométrique.
\begin{exemple}
$(u_n)$ géométrique, $u_2=12$, $q=3$. Alors $u_7=u_2\,q^{\,7-2}=12\times 3^5=12\times 243=2916$.
\end{exemple}

\methode{Déterminer la raison et le premier terme à partir de deux termes}
Traduire les deux informations en équations avec la formule du terme général, puis
résoudre le système.
\begin{exemple}
$(u_n)$ arithmétique avec $u_3=11$ et $u_8=26$. On a $u_8-u_3=(8-3)r$, soit $15=5r$,
donc $r=3$ ; puis $u_0=u_3-3r=11-9=2$.
\end{exemple}

\methode{Calculer une somme de termes consécutifs}
Repérer le nombre de termes et les termes extrêmes, puis appliquer la bonne formule.
Pour l'arithmétique : (nombre de termes)~$\times$~(moyenne des extrêmes).
\begin{exemple}
Somme $S=5+8+11+\dots+50$ (arithmétique, $r=3$). Le nombre de termes est
$\dfrac{50-5}{3}+1=16$, d'où $S=16\times\dfrac{5+50}{2}=16\times 27{,}5=440$.
\end{exemple}

\methode{Modéliser un placement financier}
Identifier le type : intérêts \emph{simples} $\to$ suite arithmétique $C_n=C_0(1+nt)$ ;
intérêts \emph{composés} $\to$ suite géométrique $C_n=C_0(1+t)^n$ ; puis appliquer.
\begin{exemple}
$300\,000$ FCFA placés à intérêts composés au taux $4\%$ pendant $3$ ans donnent
$C_3=300\,000\times 1{,}04^{\,3}\approx 300\,000\times 1{,}1249\approx 337\,459$ FCFA.
\end{exemple}
